\FSabschnitt{Teil B \quad Halbleiter dotiert $\cdot$ Fermi $\cdot$ pn-Übergang}{Klausurteil A2 / B1}



% =====================================================================
\begin{RezeptA}{3.1\ \ REZEPT: dotierter Halbleiter \hfill
  \Quelle{Ü2 A1 · Kl09 · Kl10 · Kl11-1 · Kl11-2 · PK13 A2 (6×)}}
\Gesucht{Typ (n/p), $n_0$ und $p_0$, Fermi-Abstand $W_F-W_i$,
Banddiagramm, ggf.\ $\sigma$.}

\Erkennung{„Ein Siliziumkristall ist mit $N_D=\dots$ und $N_A=\dots$
dotiert.“ $\cdot$ „Berechnen Sie die Lage des Fermi-Niveaus.“
$\cdot$ „Zeichnen Sie das Bänderdiagramm.“}

\Kette
\begin{enumerate}
\item \textbf{Typ:} $N_D>N_A\Rightarrow$ n\quad\ $N_A>N_D\Rightarrow$ p
\item \textbf{Majoritäten:} $n_0\approx N_D-N_A$ \ (n-HL),
      $p_0\approx N_A-N_D$ \ (p-HL)\\
      \textbf{Minoritäten} über $n_0\cdot p_0=n_i^2$
\item \textbf{Fermi-Abstand.} Ansatz:
      \[ n_0=n_i\,\ee^{(W_F-W_i)/k_BT}
         \notat{$\ln(\ )$, dann $\times k_BT$} \]
      \vspace{-3.2mm}
      \begin{align*}
      \text{n-HL:}\quad W_F-W_i &= k_B T\,\ln\frac{n_0}{n_i}\quad[\mathrm{J}]\\[-0.3mm]
      \text{p-HL:}\quad W_i-W_F &= k_B T\,\ln\frac{p_0}{n_i}\quad[\mathrm{J}]
      \end{align*}
      \vspace{-3.2mm}
\item \textbf{J $\to$ eV:} $\ \erg{\div\ 1{,}6\cdot10^{-19}}$
\item \textbf{Banddiagramm} zeichnen (Anleitung unten)
\end{enumerate}

\FSline
\Beispiel[$N_D=5\cdot10^{17}$, $N_A=5\cdot10^{15}\,\mathrm{cm^{-3}}$]
\textbf{1)} $5\cdot10^{17}>5\cdot10^{15}\ \Rightarrow\ \erg{\text{n-Halbleiter}}$
\vspace{0.3mm}

\textbf{2)} $n_0\approx N_D-N_A=5\cdot10^{17}-5\cdot10^{15}
 \approx\erg{5\cdot10^{17}\,\mathrm{cm^{-3}}}$\\
$p_0=\dfrac{n_i^2}{n_0}=\dfrac{(1{,}5\cdot10^{10})^2}{5\cdot10^{17}}
 =\erg{4{,}5\cdot10^{2}\,\mathrm{cm^{-3}}}$
\vspace{0.3mm}

\textbf{3)}
\vspace{-1mm}
\[\begin{aligned}
W_F-W_i &= k_B T\,\ln\frac{n_0}{n_i}\\
 &= 1{,}38\!\cdot\!10^{-23}\tfrac{\mathrm{J}}{\mathrm{K}}\cdot300\,\mathrm{K}
    \cdot\ln\frac{5\cdot10^{17}}{1{,}5\cdot10^{10}}\\
 &= 4{,}14\cdot10^{-21}\,\mathrm{J}\cdot\ln\left(3{,}33\cdot10^{7}\right)
    \notat{$k_BT$ ausrechnen}\\
 &= 4{,}14\cdot10^{-21}\,\mathrm{J}\cdot17{,}322
    \notat{$\ln(\ )$ anwenden}\\
 &= 7{,}17\cdot10^{-20}\,\mathrm{J}
    \notat{$\div\,1{,}6\cdot10^{-19}$}\\
 &= \erg{0{,}45\ \mathrm{eV}}
\end{aligned}\]

\tr{$k_BT$ \emph{zuerst} ausrechnen ($=4{,}14\cdot10^{-21}$\,J), dann erst
den $\ln$. Zehnerpotenzen mit \textsf{EXP}. Der $\ln$ ist einheitenlos —
die Einheit steckt komplett in $k_BT$.}

\FSline
\textbf{5) Banddiagramm}
\begin{center}
\begin{tikzpicture}[font=\tiny, scale=0.82]
  % Baender als Bloecke
  \fill[black!12] (0,1.60) rectangle (3.30,2.00);
  \fill[black!12] (0,-0.40) rectangle (3.30,0.00);
  \draw[thick] (0,1.60) -- (3.30,1.60);   % W_C
  \draw[thick] (0,0.00) -- (3.30,0.00);   % W_V
  \draw (0,2.00) -- (3.30,2.00);
  \draw (0,-0.40) -- (3.30,-0.40);
  \draw (0,-0.40) -- (0,2.00);
  \draw (3.30,-0.40) -- (3.30,2.00);
  % W_i und W_F
  \draw[dashed] (0,0.80) -- (3.55,0.80);
  \draw[dashed, FSrot, thick] (0,1.45) -- (3.55,1.45);
  % Beschriftung links
  \node[left=-0.3mm] at (0,1.68) {$W_C$};
  \node[left=-0.3mm] at (0,1.43) {\textcolor{FSrot}{$W_F$}};
  \node[left] at (0,0.80) {$W_i$};
  \node[left] at (0,0.00) {$W_V$};
  \node at (1.65,1.80) {Leitungsband};
  \node at (1.65,-0.20) {Valenzband};
  % Masspfeile
  \draw[<->, FSrot] (3.15,0.80) -- (3.15,1.45)
      node[midway, right=0.3mm] {$0{,}45$\,eV};
  \draw[<->] (3.00,0.00) -- (3.00,1.60)
      node[midway, right=0.3mm] {$W_g$};
\end{tikzpicture}
\end{center}
\vspace{-1mm}

\textbf{Wie man es zeichnet}
\begin{enumerate}
\item $W_V$ unten, $W_C$ oben — Abstand $=W_g=1{,}1$\,eV (Si)
\item $W_i$ \emph{mittig} zwischen $W_V$ und $W_C$
\item n-HL: \erg{$W_F$ über $W_i$} \quad p-HL: \erg{$W_F$ unter $W_i$}
\item berechneten Abstand $W_F-W_i$ als Maßpfeil eintragen
\end{enumerate}

\Variante{p-Halbleiter ($N_A>N_D$): $p_0\approx N_A-N_D$, Formel
$W_i-W_F=k_BT\ln(p_0/n_i)$, $W_F$ liegt \emph{unter} $W_i$.
Rest der Kette identisch.}
\end{RezeptA}

% =====================================================================
\begin{RezeptA}{3.2\ \ Leitfähigkeit $\sigma$ \hfill \Quelle{Ü2 A1 · Kl09 · Kl11 · PK13 A2 (4×)}}
\Gesucht{$\sigma$ bzw.\ Vergleich dotiert / undotiert.}
\Erkennung{„Wie groß ist die Leitfähigkeit?“ $\cdot$ „Um welchen Faktor
ändert sich $\sigma$ durch die Dotierung?“}

\textbf{Allgemein} \ $\sigma=e\,(\mu_n n+\mu_p p)$ \quad
$[\sigma]=\mathrm{S/m}$
\vspace{0.4mm}

\warn{Erst $\mathrm{cm^{-3}}\to\mathrm{m^{-3}}$ umrechnen ($\times10^{6}$),
sonst stimmt die Größenordnung nicht.}

\Beispiel[Si, $n_0=5\cdot10^{17}\,\mathrm{cm^{-3}}=5\cdot10^{23}\,\mathrm{m^{-3}}$]
\textbf{dotiert} — Minoritäten vernachlässigbar:
\begin{align*}
\sigma &\approx e\,\mu_n n_0\\
 &= 1{,}6\!\cdot\!10^{-19}\mathrm{As}\cdot0{,}135\tfrac{\mathrm{m^2}}{\mathrm{Vs}}
    \cdot5\!\cdot\!10^{23}\tfrac{1}{\mathrm{m^3}}\\
 &= \erg{1{,}08\cdot10^{4}\ \mathrm{S/m}}
\end{align*}
\textbf{undotiert} — $n=p=n_i=1{,}5\cdot10^{16}\,\mathrm{m^{-3}}$:
\begin{align*}
\sigma &= e\,(\mu_n+\mu_p)\,n_i\\
 &= 1{,}6\!\cdot\!10^{-19}\cdot0{,}183\cdot1{,}5\!\cdot\!10^{16}\\
 &= \erg{4{,}39\cdot10^{-4}\ \mathrm{S/m}}
\end{align*}
Verhältnis: $\dfrac{1{,}08\cdot10^{4}}{4{,}39\cdot10^{-4}}
 \approx\erg{2{,}5\cdot10^{7}}$
\tr{$\mu_n=0{,}135$ und $\mu_n+\mu_p=0{,}183\ \mathrm{m^2/Vs}$ sind
Merkwerte — stehen auf S1.}
\end{RezeptA}


% =====================================================================
\begin{RezeptA}{3.3\ \ pn-Übergang: $\rho(x)$, $E(x)$, $U(x)$ \hfill
  \Quelle{Ü2 A2 · Ü3 A1.1 · Kl09 · Kl10 · Kl12 · PK-I B1 (6×)}}
\Gesucht{Qualitativer Verlauf von Raumladung, Feldstärke und Potential
über die Raumladungszone.}
\Erkennung{„Skizzieren Sie den Verlauf von $E(x)$ und $U(x)$.“
$\cdot$ „Ankreuzen: welcher Verlauf gehört zu …“}

\textbf{Merkkette} \ $\rho(x)\ \xrightarrow{\ \int\ }\ E(x)
 \ \xrightarrow{\ \int\ }\ U(x)$ \ — jedes Integrieren macht den
Verlauf \erg{eine Stufe glatter}: Rechteck $\to$ Dreieck $\to$ S-Kurve.

\begin{center}
\begin{tikzpicture}[font=\tiny, x=0.86cm, y=0.86cm]
  % ---------- rho(x) ----------
  \begin{scope}[yshift=0cm]
    \draw[->] (-1.5,0) -- (1.7,0) node[right] {$x$};
    \draw[->] (0,-0.45) -- (0,0.55) node[above] {$\rho$};
    \fill[black!25] (-1.0,0) rectangle (0,-0.35);
    \fill[black!25] (0,0) rectangle (0.7,0.42);
    \draw[thick] (-1.5,0) -- (-1.0,0) -- (-1.0,-0.35) -- (0,-0.35)
                 -- (0,0.42) -- (0.7,0.42) -- (0.7,0) -- (1.5,0);
    \node at (-1.25,0.28) {p};
    \node at (1.15,0.28) {n};
    \node[below] at (-1.0,-0.38) {$-x_p$};
    \node[below right=-0.5mm] at (0.7,0) {$x_n$};
  \end{scope}
  % ---------- E(x) ----------
  \begin{scope}[yshift=-1.5cm]
    \draw[->] (-1.5,0) -- (1.7,0) node[right] {$x$};
    \draw[->] (0,-0.65) -- (0,0.35) node[above] {$E$};
    \draw[thick, FSrot] (-1.5,0) -- (-1.0,0) -- (0,-0.55) -- (0.7,0) -- (1.5,0);
    \node[FSrot, right=0.5mm] at (0.05,-0.5) {$E_{max}$};
    \node[below] at (-1.0,0) {$-x_p$};
  \end{scope}
  % ---------- U(x) ----------
  \begin{scope}[yshift=-3.0cm]
    \draw[->] (-1.5,0) -- (1.7,0) node[right] {$x$};
    \draw[->] (0,-0.15) -- (0,0.75) node[above] {$U$};
    \draw[thick] (-1.5,0) -- (-1.0,0)
      .. controls (-0.5,0.02) and (-0.35,0.28) .. (0,0.30)
      .. controls (0.35,0.32) and (0.55,0.55) .. (0.7,0.57) -- (1.5,0.57);
    \draw[<->, FSrot] (1.25,0) -- (1.25,0.57)
        node[midway, right=0.3mm] {$U_D$};
  \end{scope}
\end{tikzpicture}
\end{center}
\vspace{-1.5mm}

\textbf{Merksätze}
\begin{itemize}
\item $\rho$: p-Seite \erg{negativ} (ortsfeste Akzeptor-Rümpfe),
      n-Seite \erg{positiv} (Donator-Rümpfe)
\item $E(x)$ dreieckig, \erg{Betragsmaximum bei $x=0$}, null an den
      RLZ-Rändern
\item $U(x)$ s-förmig, Gesamthöhe $=U_D$
\item Ladungsneutralität: $\ N_A\,x_p=N_D\,x_n$ \ $\Rightarrow$
      schwächer dotierte Seite hat die \erg{breitere} RLZ
\item \textbf{Diffusion} treibt Träger ins Gegengebiet,
      \textbf{Drift} im E-Feld treibt sie zurück — im Gleichgewicht
      heben sie sich auf
\end{itemize}
\end{RezeptA}

% =====================================================================
\begin{RezeptA}{3.4\ \ Diffusionsspannung $U_D$ \hfill \Quelle{Ü3 A1 · Kl10 · Kl12 (3×)}}
\Gesucht{Eingebaute Spannung des unbelasteten pn-Übergangs.}
\Erkennung{„Berechnen Sie die Diffusionsspannung.“ $\cdot$ gegeben sind
$N_A$, $N_D$ und $n_i$.}
\[ U_D=\frac{k_B T}{e}\,\ln\frac{N_A\,N_D}{n_i^{2}}
   = \UT\,\ln\frac{N_A N_D}{n_i^{2}} \]

\Beispiel[$N_A=N_D=10^{15}\,\mathrm{cm^{-3}}$]
\begin{align*}
U_D &= 25{,}9\,\mathrm{mV}\cdot\ln
       \frac{10^{15}\cdot10^{15}}{(1{,}5\cdot10^{10})^{2}}\\
 &= 25{,}9\,\mathrm{mV}\cdot\ln\left(4{,}44\cdot10^{9}\right)
    \notat{Bruch zuerst}\\
 &= 25{,}9\,\mathrm{mV}\cdot22{,}21
    \notat{$\ln(\ )$ anwenden}\\
 &= \erg{575\ \mathrm{mV}}
\end{align*}
\warn{Mit $\UT=25$\,mV kämen 555\,mV heraus.} Den in der Aufgabe
vorgegebenen Wert nehmen — s.\ S1.
\tr{$n_i^2$ nicht vergessen zu quadrieren: $(1{,}5\cdot10^{10})^2
=2{,}25\cdot10^{20}$.}
\end{RezeptA}

% =====================================================================
\begin{RezeptB}{3.5\ \ Potentialtopf \hfill \Quelle{Ü1 A1 · Kl09 · Kl10 · PK-I A1 (4×)}}
\Erkennung{„unendlich hoher Potentialtopf“ $\cdot$ „Breite $L$“
$\cdot$ „Sind die Energieeigenwerte gequantelt oder kontinuierlich?“}
\[ E_n=\frac{n^{2}\pi^{2}\hbar^{2}}{2\,m\,L^{2}}\qquad
   \Delta E=E_2-E_1=\frac{3\pi^{2}\hbar^{2}}{2mL^{2}} \]
$\hbar=h/2\pi$ $\cdot$ Ergebnis in J, dann
\erg{$\div\,1{,}6\cdot10^{-19}$} für eV\\
\warn{$m=10^{-30}$\,kg} ist Klausurvorgabe (Ü01: $9{,}1\cdot10^{-31}$\,kg)
\vspace{0.4mm}

\textbf{Antwort auf die Ankreuzfrage:} Energieeigenwerte sind
\erg{gequantelt}. $L$ größer $\Rightarrow$ $\Delta E$ kleiner
$\Rightarrow$ quasi-kontinuierlich (Übergang zum Festkörper).
\end{RezeptB}

% =====================================================================
\begin{RezeptB}{3.6\ \ Photon \& Materiewelle \hfill \Quelle{Kl10 A24/25 · PK13 A1 (2×)}}
\Erkennung{„Welche Wellenlänge …“ $\cdot$ „innerer Fotoeffekt“
$\cdot$ „de-Broglie-Wellenlänge“}
Photonenenergie: $W=h\,f=\dfrac{h\,c}{\lambda}$\\
innerer Fotoeffekt: $\lambda_{max}=\dfrac{h\,c}{W_g}$\qquad
de Broglie: $\lambda=\dfrac{h}{m\,v}$
\vspace{0.3mm}

Si: $\lambda_{max}=\dfrac{6{,}6\cdot10^{-34}\cdot3\cdot10^{8}}
{1{,}1\cdot1{,}6\cdot10^{-19}}=\erg{1{,}13\ \mu\mathrm{m}}$ (IR)
\vspace{0.3mm}

\textbf{Merksatz:} nur $\lambda<\lambda_{max}$ erzeugt Paare —
größere Wellenlänge $=$ zu wenig Energie.
\tr{$W_g$ in eV $\to$ J: $\times\,1{,}6\cdot10^{-19}$ \emph{vor} dem Teilen.}
\end{RezeptB}

% =====================================================================
\begin{RezeptC}{3.7\ \ Poisson / Raumladung \hfill \Quelle{Ü2 A2 · Kl09/10/12 qualitativ}}
$\dfrac{\dd^{2}\varphi}{\dd x^{2}}=-\dfrac{\rho}{\eps}$ \quad
$E(x)=\displaystyle\int\frac{\rho}{\eps}\,\dd x$ \quad
$U=-\displaystyle\int E\,\dd x$\\
Ladungsneutralität: $N_A x_p=N_D x_n$;\ \ RLZ-Weite
$w=x_p+x_n$
\end{RezeptC}


